\newpage\pagestyle{conclusions}

\unnumberedchapter{Conclusioni e Sviluppi Futuri}

\section*{Sintesi dei risultati ottenuti}
Il presente lavoro di tesi ha esplorato l'integrazione di database a grafo distribuiti all'interno di architetture di \textit{Conversational AI}, con l'obiettivo di superare i limiti di memoria e latenza tipici dei sistemi RAG tradizionali. Attraverso l'implementazione di una struttura a due livelli (Remote e Local Shard), è stato possibile creare un agente virtuale in Unreal Engine capace di attingere a una base di conoscenza vasta senza compromettere la fluidità dell'interazione.

I test prestazionali condotti nel Capitolo 5 hanno confermato l'efficacia della logica di sharding: l'uso di una cache locale basata su algoritmi di \textit{eviction} LFU ha permesso di abbattere i tempi di risposta fino al 68\% rispetto a un'architettura monolitica. Inoltre, l'ottimizzazione tramite query parametrizzate ha dimostrato come una corretta gestione dell' \textit{Execution Plan Cache} in Neo4j sia fondamentale per garantire scalabilità e stabilità in ambienti di produzione.

\section*{Limitazioni attuali}
Nonostante i risultati positivi, il sistema presenta alcune limitazioni che aprono la strada a ulteriori ricerche. Attualmente, la sincronizzazione tra il database remoto e lo shard locale avviene principalmente in modalità \textit{pull}: lo shard viene aggiornato solo a seguito di una richiesta dell'utente. Questo implica che, se un documento venisse modificato sul database remoto, lo shard locale potrebbe servire informazioni non aggiornate fino alla successiva operazione di sincronizzazione o fino alla scadenza del nodo nella cache.

\newpage

\section*{Sviluppi Futuri}
Le direzioni future per l'evoluzione di questo progetto si concentrano sulla resilienza del dato e sulla scalabilità dell'infrastruttura:

\begin{itemize}
    \item \textbf{Integrazione con Apache Kafka}: Per risolvere il problema della coerenza del dato tra Remote e Shard, una naturale evoluzione prevede l'impiego di \textbf{Apache Kafka} \cite{apache_kafka}. Implementando un paradigma \textit{Event-Driven} \cite{shapira2021kafka}, ogni aggiornamento effettuato sulla base di conoscenza remota verrebbe pubblicato su un topic Kafka. Gli shard locali, agendo come \textit{consumer}, potrebbero aggiornare istantaneamente i propri nodi in tempo reale, garantendo che l'agente virtuale disponga sempre dell'ultima versione disponibile della documentazione.

    \item \textbf{Orchestrazione tramite Kubernetes}: In ottica di distribuzione su larga scala, il sistema potrebbe beneficiare dell'orchestrazione tramite \textbf{Kubernetes} \cite{kubernetes}. L'intera architettura (server LangGraph, istanze Neo4j e backend Python) potrebbe essere containerizzata, permettendo una gestione dinamica delle risorse. Kubernetes consentirebbe di scalare orizzontalmente il numero di shard locali in base al numero di utenti attivi o alla distribuzione geografica, garantendo alta affidabilità e bilanciamento del carico.
\end{itemize}